\section{Chebotarev: an exact target and a proved effective reduction}
\label{sec:chebotarev}

\subsection{The theorem that remains to be internalized}

Let $L/K$ be a finite Galois extension of number fields, let
$G=\operatorname{Gal}(L/K)$, and let $C$ be a conjugacy class in $G$.
For a nonzero prime ideal $\mathfrak p$ of $\mathcal O_K$ unramified in $L$,
the Frobenius class is the conjugacy class of the automorphism determined
modulo a prime $\mathfrak P$ above $\mathfrak p$ by
\[
 \operatorname{Frob}_{\mathfrak P}(z)\equiv
 z^{N\mathfrak p}\pmod{\mathfrak P}.
\]
Set
\[
 \Pi_C(x)=\#\{\mathfrak p:N\mathfrak p\leq x,\
      \mathfrak p\text{ unramified},\ \operatorname{Frob}_{\mathfrak p}=C\},
 \qquad
 \Pi_K(x)=\#\{\mathfrak p:N\mathfrak p\leq x\}.
\]
The natural-density conclusion is
\[
 \frac{\Pi_C(x)}{\Pi_K(x)}\longrightarrow\frac{|C|}{|G|}.
\]
The customary asymptotic $\Pi_C(x)\sim(|C|/|G|)\Li(x)$ implies it together
with the prime ideal theorem. No density conclusion follows merely from
the ability to enumerate the Frobenius classes below a finite bound.

Put $g=|G|$ and $c=|C|$. Once the number-field input has been correctly
coded, a completely integer target is
\begin{equation}
 \forall k\geq1\ \exists M\ \forall x\geq M:
 \quad \Pi_K(x)>0\ \land\
 k\,|g\Pi_C(x)-c\Pi_K(x)|\leq g\Pi_K(x).
 \label{eq:chebotarev-integer-target}
\end{equation}
Its constructive version should provide $M$ as a specified primitive
recursive function of the field data and $k$, with a proof of its matrix
in $\ipra$. The existence of classical asymptotic notation alone does not
supply that object.

\subsection{A finite representation of the input}

Here is a concrete representation to use in a formalization. It is a
specification for the required algebraic layer, not an assertion that this
layer has already been implemented.

Represent each number field by a monic irreducible polynomial over $\Q$,
the rational power-basis arithmetic, and a certified integral basis. An
embedding $K\hookrightarrow L$ is a rational matrix preserving $1$ and
multiplication. Represent $G$ by a list of rational matrices with a group
multiplication table and certificates that these are all $K$-automorphisms
of $L$; in a Galois presentation the list has length $[L:K]$. Represent $C$
by a finite sublist and check conjugacy using this multiplication table.

For a degree-$d$ field with a certified integral basis, an integral ideal
of norm at most $x$ is represented by a Hermite normal form matrix. Its
positive diagonal entries have product at most $x$, and the usual reduced
off-diagonal entries are bounded by diagonal entries, hence by $x$.
Thus enumeration over matrices with entries bounded by $x$ includes every
such ideal. Multiplicative stability under the integral basis is a finite
integer divisibility check. The quotient is a finite ring of at most $x$
elements. A proper ideal is prime exactly when its quotient is a field;
checking that every nonzero residue has an inverse is a bounded check.
Norm is the determinant, and equality is equality of the canonical matrix.

To inspect a prime above $\mathfrak p$ one may enumerate ideals of
$\mathcal O_L$ of norm at most $(N\mathfrak p)^{[L:K]}$. The algebraic
theorem ensuring that a prime above it occurs in this list is part of the
required library. Ramification, decomposition and the Frobenius test then
reduce to finite ideal and quotient calculations. The conjugacy-class
answer is independent of the chosen prime above $\mathfrak p$ by the
standard finite Galois argument, which must also be proved in the encoded
setting.

Certification of the full ring of integers, complete automorphism list,
factorization of ideals, and these independence statements cannot simply
be replaced by unchecked Boolean labels. Either they are established from
the field presentation with explicit resource bounds, or the resulting
theorem is explicitly conditional on supplied certificates whose meanings
have themselves been proved. The present package chooses the second
interface for its targets and does not discharge those obligations.

\subsection{Effectivity does not fail merely because of a possible zero}

Unconditional effective Chebotarev estimates exist. Their typical structure
is a main term, a possible real exceptional-zero contribution, and a
controlled decaying error. Effective bounds on the distance of that zero
from $1$ in terms of field data then give an effective density modulus
\cite{LagariasOdlyzko,ThornerZaman2019}. The much stronger uniform
conductor dependence needed for Siegel--Walfisz is a different issue.

The following proposition proves the numerical part of that passage
completely. It is deliberately stated with an explicit analytic input, so
that its conclusion cannot be mistaken for an independent proof of
Chebotarev.

\begin{proposition}[A primitive recursive density modulus from supplied estimates]
\label{prop:chebotarev-modulus}
Let $n,k\geq1$, $b\geq1$, $c_0,d,U_0\in\N$, and $0\leq\rho\leq1$.
Suppose counting functions $P,Q$ satisfy, for $x\geq e^{U_0}$ and $x\geq4$,
\begin{align*}
 |P(x)-\rho\Li(x)|&\leq R(x),\\
 |Q(x)-\Li(x)|&\leq R(x),
\end{align*}
where
\[
 R(x)=2x\exp(-2^{-b}\log x)
       +2^{c_0}x\exp\!\left(-2^{-d}\sqrt{\frac{\log x}{n}}\right).
\]
Define
\begin{align*}
 P_0&=2\left(U_0+64k\,2^{2b}
             +384k\,2^{c_0+4d}n^2+4\right),\\
 X_0&=2^{P_0}.
\end{align*}
For every integer $x\geq X_0$ one has $Q(x)>0$ and
\[
 \left|\frac{P(x)}{Q(x)}-\rho\right|\leq\frac1k.
\]
If $P=\Pi_C$, $Q=\Pi_K$, and $\rho=c/g$, this gives the matrix of
\eqref{eq:chebotarev-integer-target} with $M=X_0$.
\end{proposition}

\begin{proof}
Put $u=\log x$. Since $\log2\geq1/2$, the assumption on $x$ gives
$u\geq P_0/2$, so every assumed estimate applies. Integrating just from
$\sqrt x$ to $x$ gives, for $x\geq4$,
\[
 \Li(x)\geq\frac{x-\sqrt x}{\log x}\geq\frac{x}{2u}.
\]
For $v>0$, positivity of the exponential series gives
$e^v\geq v^2/2$ and $e^v\geq v^4/24$. Consequently
\begin{align*}
 \frac{2x e^{-2^{-b}u}}{\Li(x)}
 &\leq4u e^{-2^{-b}u}
 \leq\frac{8\,2^{2b}}u\leq\frac1{8k},\\
 \frac{2^{c_0}x e^{-2^{-d}\sqrt{u/n}}}{\Li(x)}
 &\leq 2^{c_0+1}u e^{-2^{-d}\sqrt{u/n}}
 \leq\frac{48\,2^{c_0+4d}n^2}u\leq\frac1{8k}.
\end{align*}
Thus $R(x)\leq\eta\Li(x)$ with $\eta=1/(4k)\leq1/4$.
It follows that $Q(x)\geq(1-\eta)\Li(x)>0$, and
\[
 \left|\frac{P(x)}{Q(x)}-\rho\right|
 \leq\frac{|P(x)-\rho\Li(x)|+\rho|\Li(x)-Q(x)|}{Q(x)}
 \leq\frac{2\eta}{1-\eta}\leq\frac1k.
\]
Every chosen threshold is a composition of integer sums, products and
exponentials. The same computation with rational error allowances yields
the integer conclusion by the rational-gap lemma.
\end{proof}

For example, an exceptional contribution bounded by
$\Li(x^\beta)$ with $1/2\leq\beta\leq1-2^{-b}$ is covered by the first
term in $R$: for $x\geq4$, $x^\beta\geq2$ and
\[
 \Li(x^\beta)\leq 2x^\beta
 \leq2x\exp(-2^{-b}\log x).
\]
The first inequality uses $1/\log t\leq2$ for $t\geq2$.
When no exceptional term occurs, the same bound remains valid by adding
a nonnegative allowance. Thus the reduction need not decide the existence
of an exceptional zero. It needs a proved estimate valid in both cases.

\subsection{The remaining mathematical work}

To turn Proposition~\ref{prop:chebotarev-modulus} into the requested
unconditional $\ipra$ proof, one must supply all of the following with
arithmetic derivations:
\begin{enumerate}
\item The certified number-field and prime-ideal algorithms described
above, including the correctness of Frobenius and its conjugacy class.
\item Quantitative analytic continuation, functional equations and
explicit-formula estimates for the necessary Dedekind and Hecke
$L$-functions, using finite contours with separation data.
\item A zero-free estimate and an effective lower separation for any
exceptional zero, with the dependence on the field presentation tracked.
\item Primitive recursive integers $b,c_0,d,U_0$ and proofs that they
give the two estimates required in the proposition, including the prime
ideal main term for $K$ and the finitely many ramified primes.
\end{enumerate}

The cited effective theorems show that the classical analytic program is
available. They do not certify its use of induction, the field-construction
algorithms, or its formal translation to $\ipra$. These are substantial
unfilled obligations in this package. No implication from ``effective'' to
``proved in $\ipra$'' is being used.

For a cyclotomic extension of $\Q$, Frobenius sends a root of unity to its
$p$-th power, so its conjugacy condition is the familiar residue-class
condition. Our direct Dirichlet theorem therefore supplies the corresponding
\emph{infinitude} consequence. It does not supply the natural-density
statement, even in this special case.
